\section{Introduction}\label{sec:intro}

BIKE and other QC-MDPC designs rest on iterative bit-flipping (BF) decoding, whose rare failures shape their security arguments~\cite{misoczki2013mdpc,bike2024spec,alagic2025nist}. The Level-1 target is a decoding failure rate (DFR) near $2^{-128}$, far below direct simulation~\cite{sendrier2019dfr,vasseur2021thesis}. Failures are not spread evenly: they cluster on particular keys~\cite{guo2016key,wang2023weak,nosouhi2023weakkey}, and an adversary who can detect or time them recovers the key~\cite{eaton2018timing}. A decoder's per-round thresholds form a schedule that can be retuned without touching key generation. We therefore test whether a deliberate change of schedule alters failure risk for a given key.

While characterising the CLASSIC BF decoder of Vasseur's reference implementation~\cite{vasseur2021thesis}, which we denote C0, we noticed that its earliest decision is unusually delicate. Round one acts on the untouched syndrome, so a mistake there enters the state of every later round. Raising the round-one threshold by two lowers the conditional failure rate by 4.15 in the random condition and 5.47 under an adversarial probe; lowering it by one raises the rate about eightfold. A two-unit offset is also destructive when sustained from round two onward, yet helpful when confined to one round.

We rebuilt the experiment around paired instances. An instance is a key $H$, an error $e$ and its syndrome $s=He^{\top}$; every schedule variant decodes the same triple. FRO-$k$ raises the round-one threshold by $k$ in round one only, and FRONEG lowers it by one. We summarise comparisons by the risk ratio (RR), the baseline failure rate divided by the variant's, so values above one favour the variant. Because failures concentrate on keys, intervals and tests resample keys, not errors. All quoted rates lie where failures are frequent enough to count; we make no claim about the production tail.

The study has four contributions. First, the first round is a two-sided window. Conservative offsets of $+1$ to $+3$ lower the failure rate 2.3--5.5-fold, while the permissive side collapses within two counts (Section~\ref{sec:formal}). Second, the same offset changes sign with its round and persistence. The $+2$ applied once in round two gives a small benefit, but sustained from round two onward it gives RR 0.0168 (Section~\ref{sec:temporal}). Third, the mechanism is asymmetric. FRO-2 reduces incorrect first-round flips about tenfold while correct flips fall far less (Section~\ref{sec:mechanism}). Fourth, the window reproduces at the production block size once the error weight is raised, while at the production weight only a decision margin and a convergence cost are measurable (Section~\ref{sec:production}). A designed negative closes the list: a round-one surrogate did not predict the effect of another schedule (Section~\ref{sec:counterfactual}).

Section~\ref{sec:background} sets up QC-MDPC BF decoding, related work, and a first-order model of the round-one window. Section~\ref{sec:method} gives the paired design, and Section~\ref{sec:results} the results. Section~\ref{sec:discussion} interprets them and states their limits; Section~\ref{sec:conclusion} concludes.
